% A made-up search paper for the Kovrune physics suite (eval/physics/README.md). The collaboration,
% detector and results are invented. Written as particle-physics papers on arXiv often are: hepunits for
% units, the paper's own macros for luminosity and particle names, numbers with asymmetric uncertainties.
\documentclass[aps,prd,twocolumn,nofootinbib]{revtex4-2}
\usepackage{amsmath}
\usepackage{hepunits}
\usepackage{graphicx}
\newcommand{\ifb}{\ensuremath{\text{fb}^{-1}}}
\newcommand{\Zprime}{\ensuremath{Z'}}
\newcommand{\ZSSM}{\ensuremath{Z'_\mathrm{SSM}}}
\newcommand{\Zpsi}{\ensuremath{Z'_\psi}}
\newcommand{\pt}{\ensuremath{p_\mathrm{T}}}
\newcommand{\mll}{\ensuremath{m_{\ell\ell}}}
\newcommand{\sqs}{\ensuremath{\sqrt{s}}}

\begin{document}

\title{Search for high-mass resonances in dielectron and dimuon final states with \mbox{139 \ifb} of proton--proton collisions at \sqs{} = 13\,\TeV{} with the HORIZON detector}

\author{The HORIZON Collaboration}

\begin{abstract}
A search for new resonances decaying into two electrons or two muons is presented, using \mbox{139\,\ifb} of proton--proton collisions at \sqs{} = 13\,\TeV{} recorded by the HORIZON detector from 2015 to 2018. The dilepton invariant mass spectrum is examined between 250\,\GeV{} and 6\,\TeV. No significant excess over the Standard Model expectation is observed. The largest deviation, in the dielectron channel at a mass of 774\,\GeV, has a local significance of $3.1\sigma$ and a global significance of $1.4\sigma$. Upper limits are set on the cross section times branching ratio of a new resonance. The sequential Standard Model \ZSSM{} is excluded below a mass of 5.1\,\TeV{} at 95\% confidence level (CL), and the \Zpsi{} of the $E_6$ model below 4.5\,\TeV.
\end{abstract}

\maketitle

\section{Introduction}

Many extensions of the Standard Model (SM) predict new heavy neutral gauge bosons, generically called \Zprime{} bosons, that couple to leptons. A resonance decaying into an electron or muon pair would appear as a narrow peak in the dilepton invariant mass \mll{} above a smoothly falling background, dominated by the Drell--Yan process. This makes the dilepton channels among the most sensitive to such states.

This paper uses the full data set of the second running period, with an integrated luminosity of 139\,\ifb{} at \sqs{} = 13\,\TeV. Two benchmark models are considered: the sequential Standard Model (SSM), in which the \ZSSM{} has the same couplings to fermions as the SM $Z$ boson, and the \Zpsi{} of the $E_6$ grand unified theory, whose couplings are weaker.

\section{Event selection}

Events were recorded with dielectron and single-muon triggers. Electron candidates must have $\pt > 30\,\GeV$ and pseudorapidity $|\eta| < 2.47$, excluding the transition region $1.37 < |\eta| < 1.52$ between the barrel and endcap calorimeters. Muon candidates must have $\pt > 30\,\GeV$ and $|\eta| < 2.5$. Both leptons must be isolated, and events with two same-flavour leptons are kept; the two channels do not overlap.

For a \ZSSM{} with a mass of 2\,\TeV, the signal efficiency, including acceptance, trigger, reconstruction and selection, is 71\% in the dielectron channel and 48\% in the dimuon channel. The lower dimuon efficiency comes from the stringent requirements on the muon spectrometer hits that keep the momentum resolution at high \pt{} under control.

The dielectron mass resolution is 1.0\% at 1\,\TeV. The dimuon mass resolution degrades with mass and reaches 8\% at 3\,\TeV, because the curvature of high-momentum muon tracks is small.

\section{Background estimate}

The background is dominated by Drell--Yan production, with smaller contributions from top-quark pairs, single top quarks, dibosons and, in the dielectron channel, multijet events with misidentified electrons. The \mll{} distribution is fitted with a smooth function whose normalization is constrained by the Drell--Yan simulation. The fitted normalization of the Drell--Yan background relative to the simulation is $1.05^{+0.09}_{-0.08}$, consistent with unity.

\section{Systematic uncertainties}

The main systematic uncertainties on the signal and on the background are listed in Table~\ref{tab:syst}. They are evaluated at the masses where they matter most. The theoretical uncertainty from the parton distribution functions (PDF) dominates at high mass.

\begin{table}
\caption{Main systematic uncertainties on the expected signal yield, at the mass given for each source.}
\label{tab:syst}
\begin{ruledtabular}
\begin{tabular}{lcc}
Source & Mass & Uncertainty \\
\hline
Luminosity & all & 1.7\% \\
Electron energy scale & 1\,\TeV & 0.5\% \\
Muon momentum scale & 3\,\TeV & 2.0\% \\
Trigger efficiency & all & 1.0\% \\
PDF & 5\,\TeV & 8\% \\
\end{tabular}
\end{ruledtabular}
\end{table}

\section{Results}

The observed \mll{} spectra agree with the background expectation over the whole range searched. The largest local excess is found in the dielectron channel at $\mll = 774\,\GeV$, with a local significance of $3.1\sigma$. After accounting for the look-elsewhere effect over the mass range from 250\,\GeV{} to 6\,\TeV, its global significance is $1.4\sigma$. No excess above $2\sigma$ local significance is seen in the dimuon channel.

Upper limits at 95\% CL are set on the cross section times branching ratio, $\sigma B$, of a resonance decaying into a lepton pair, using a profile-likelihood test statistic and the CL$_\mathrm{s}$ method. For a resonance with a natural width of 3\% of its mass, the combined upper limit on $\sigma B$ falls from 4.5\,fb at 500\,\GeV{} to 0.020\,fb at 6\,\TeV.

From these limits, lower limits on the resonance masses are derived. The \ZSSM{} is excluded below 4.8\,\TeV{} in the dielectron channel, below 4.6\,\TeV{} in the dimuon channel, and below 5.1\,\TeV{} when the two are combined. The \Zpsi{}, whose cross section is smaller, is excluded below 4.5\,\TeV{} in the combination. Masses above these limits remain allowed.

\section{Conclusion}

A search for dilepton resonances with 139\,\ifb{} of data at \sqs{} = 13\,\TeV{} finds no significant excess. The sequential Standard Model \ZSSM{} is excluded below 5.1\,\TeV{} and the \Zpsi{} below 4.5\,\TeV, at 95\% CL. These limits improve on the previous HORIZON results with 36\,\ifb{} by about 1\,\TeV.

\begin{acknowledgments}
This paper is made up for the Kovrune physics suite; the HORIZON Collaboration does not exist.
\end{acknowledgments}

\begin{thebibliography}{2}
\bibitem{Langacker} P. Langacker, Rev. Mod. Phys. \textbf{81}, 1199 (2009).
\bibitem{Cowan} G. Cowan, K. Cranmer, E. Gross and O. Vitells, Eur. Phys. J. C \textbf{71}, 1554 (2011).
\end{thebibliography}

\end{document}
